% ===== §19 =====
\section{The Limit as Ground}
\label{sec:close}

\noindent
This section closes the paper without gathering its results into a final synthesis. Its aim is to state the change of aspect the whole argument has effected, that the limit of language, taken by a first habit for a defect, has shown itself throughout as a ground, and then to leave several voices speaking in their unresolved tension, in keeping with an argument that has warned against the closure of what should stay open. The section states the change of aspect, sets out the polyphony, and returns the thesis to the open.

\subsection*{Why no synthesis}

The paper has followed a single thesis through six dimensions, a formal core, a persistence, and three registers of generativity. It closes without gathering these into a final synthesis, and the refusal is not a modesty but a consequence of the argument. A work about what language cannot reach ought not to end by reaching it. To compose the results into a single culminating proposition would be to perform on the paper's own thesis the closure the thesis identifies as extinction, to still the movement of the argument by delivering it, finished, as a thing possessed. The form of the close must enact the content of the paper, and the content forbids the closure that a synthesis would be.

\subsection*{The change of aspect}

What can be said in closing is a change of aspect rather than a summary. The limit of language, which the first habit took for a defect, has shown itself throughout as a ground. In existence it is what keeps the bond an act rather than a title. In maintenance it is what returns the relation to presence rather than record. In value it is what keeps worth living in the bond rather than available to be drawn off. In ethics it is what returns the response to a perception no rule replaces. In justice it is what keeps the law answerable to a generativity it cannot close. In evolution it is the space in which a relation can become what no one yet has words for. Beneath all six, it is the manque around which drive circles and along whose failure desire slides, the open center without which the relation would not move. The same open center, in Part Four, was shown to be the condition of the poem, of an ethics without possession, and of a justice without premature settlement. The limit is not a problem set for relational life to solve. It is the ground on which relational life stands.

\subsection*{The voices}

Several voices may be left to speak here, and the paper does not silence them into one, for the truth of the matter lies in their tension and not in any one of them.

There is the voice that hears in all this a hard necessity, the irreducible solitude of subjects who cannot finally reach one another, and grants that the account does not dissolve that solitude but only reads it otherwise. This voice is not refuted by the paper. The subjects of a bond do remain, at the core, unreached by each other, and if one had wanted the fusion of two into a transparency without remainder, the account gives it cold comfort, telling one only that the fusion so wanted would have been a death.

There is the voice that hears a liberation, the release from a wish for total transparency that was never a wish for love but a wish for possession, and finds in the limit the room where love could breathe. To this voice the same result is good news: what one feared was a poverty, the unreachability of the beloved, is the very room in which the beloved stays a living source of love, and the letting-go of the wish to possess is not a loss but the beginning of a bond that can last.

There is the voice that hears a discipline, the long schooling of a perception and a patience fine enough to dwell with what cannot be closed. To this voice the account is neither cold nor consoling but a task: to learn the aesthetic measure that ethics asks, to practice the witnessing that justice asks, to make and keep the room that intimacy asks, none of which is given and all of which must be cultivated over a life.

These voices are not reconciled here. To compose them into a single conclusion would be to do to the argument what the argument warns against, closing a gap whose openness was the whole point. They are left in their tension, three ways of standing toward one result, and the paper does not adjudicate among them.

\subsection*{The thesis returned to the open}

So the paper ends by returning its thesis to the open. The limit of language makes drive possible, makes the poem possible, makes an ethics without possession and a justice without premature settlement possible, and makes possible a bond that keeps moving because it never closes. What language cannot reach is not the failure of relational life. It is the opening in which relational life begins, and goes on, and does not end.
